\section{Fases de desarrollo previstas y cronograma}
\label{sec:fases}

El desarrollo del presente TFG constará de diferentes fases que se describen a lo largo de esta sección.

\vspace{-0.15cm}

\subsection{Definición del problema}
En esta fase el objetivo es definir cuál es el objeto y tema sobre el cual se va a desarrollar el proyecto, en este caso, una aplicación web para el seguimiento del estado de ánimo.

\vspace{-0.15cm}

\subsection{Estudio del problema}
En esta fase se definirá claramente la cuestión a desarrollar, se hará un estudio profundo del tema y se establecerán cuáles son los objetivos del proyecto.

\vspace{-0.15cm}

\subsection{Especificación y análisis de requisitos}
En esta fase se recopilarán, analizarán y verificarán las necesidades que deberá cumplir la aplicación web \texttt{MoodNest}. Se recogerán en forma de requisitos funcionales, no funcionales y requisitos de información.

\vspace{-0.15cm}

\subsection{Diseño}
En esta fase se hará el desarrollo del diseño conceptual del sistema. Se definirá la experiencia de usuario y se crearán los prototipos de la interfaz. Además, se diseñará la arquitectura \textit{software}, el modelo de la base de datos y el enfoque de la seguridad del sistema.

\vspace{-0.15cm}

\subsection{Implementación}

En esta fase se transformará el diseño conceptual, definido en la Sección \ref{ssec:diseno}, en código ejecutable, materializando la arquitectura API REST + SPA. Esta es la fase principal de desarrollo de \textit{software} y se dividirá en cuatro grandes tareas:

\begin{itemize}
    \item \textbf{Desarrollo del \textit{Backend} (API REST):} Se codificarán los Controladores, Servicios y Repositorios en \texttt{Java} utilizando el \textit{framework} \texttt{Spring Boot}. Se implementará toda la capa de acceso a datos para gestionar la lógica de negocio y las consultas a la base de datos \texttt{MongoDB}.
    
    \item \textbf{Desarrollo del \textit{Frontend} (SPA):} Se construirán los componentes visuales de la interfaz de usuario utilizando la biblioteca \texttt{React.js}. Esto incluye la maquetación de las distintas vistas y la integración de las librerías 
    de gráficos.
    
    \item \textbf{Integración \textit{Frontend-Backend}:} Se conectará la aplicación cliente con la \texttt{API} del servidor, asegurando que el flujo de datos en formato \texttt{JSON} entre ambas partes sea correcto y robusto.

    \item \textbf{Containerización (Docker):} Se crearán los archivos \texttt{Dockerfile} para el \textit{frontend} y el textit{backend}, y un archivo \texttt{docker-compose.yml} que orquestará la aplicación completa (Frontend, Backend y Base de Datos) para su despliegue y evaluación.
\end{itemize}

Para todo este proceso se emplearán las tecnologías y herramientas específicas ya detalladas en la Sección \ref{ssec:implementacion} de este documento.

\vspace{-0.3cm}

\subsection{Pruebas y tests}
\vspace{-0.05cm}
Durante el desarrollo del proyecto se procederá a realizar las pruebas necesarias para asegurar que se cumplen con todos los requisitos especificados. Se realizarán pruebas unitarias, pruebas de integración, pruebas funcionales y pruebas de usabilidad de la interfaz, con el objetivo de encontrar todos los posibles errores y asegurar el uso óptimo de la aplicación.

\vspace{-0.3cm}

\subsection{Documentación}
\vspace{-0.05cm}
Durante todo el desarrollo del proyecto se elaborará la presente memoria, donde se detallará todo lo relativo a la aplicación web, facilitando el entendimiento y comprensión del mismo a personas que quieran información del proyecto o deseen utilizar la aplicación.

\vspace{-0.3cm}

\subsection{Cronograma inicial}
\vspace{-0.05cm}
El TFG consta de 12 créditos, lo que según la relación establecida créditos/horas corresponde a 300 horas de trabajo. 

En la Tabla \ref{tab:cronograma} se detalla una estimación de las horas asignadas a cada fase divididas en 4 meses, que es el tiempo estimado para la realización del TFG.

\begin{table}[h!]
    \centering
    \begin{tabular}{|l|c|c|c|c|c|}
        \hline
        \textbf{Tarea} & \textbf{Mes 1} & \textbf{Mes 2} & \textbf{Mes 3} & \textbf{Mes 4} & \textbf{Total} \\
        \hline
        Definición del problema & 5 & & & & \textbf{5} \\
        \hline
        Estudio del problema & 15 & & & & \textbf{15} \\
        \hline
        Especificación y análisis de requisitos & 25 & 5 & & & \textbf{30} \\
        \hline
        Diseño & 20 & 25 & & & \textbf{45} \\
        \hline
        Implementación & & 30 & 45 & 20 & \textbf{95} \\
        \hline
        Pruebas & & & 10 & 20 & \textbf{30} \\
        \hline
        Documentación & 10 & 15 & 25 & 30 & \textbf{80} \\
        \hline
        \textbf{Total} & \textbf{75} & \textbf{75} & \textbf{80} & \textbf{70} & \textbf{300} \\
        \hline
    \end{tabular}
    \caption{Cronograma del TFG.}
    \label{tab:cronograma}
\end{table}